\documentclass[11pt,a4paper]{article}
\usepackage[utf8]{inputenc}
\usepackage[T1]{fontenc}
\usepackage{lmodern}
\usepackage[french]{babel}
\usepackage{amsmath,amssymb,mathtools}
\usepackage{icomma}
\usepackage{xcolor}
\usepackage{tikz}
\usepackage{pgfplots}
\usepackage[most]{tcolorbox}
\usepackage{enumitem}
\usepackage{fancyhdr}
\usepackage[hidelinks]{hyperref}
\usepackage{titlesec}
\usepackage[a4paper,margin=2.2cm,headheight=26pt,headsep=10pt]{geometry}
\usetikzlibrary{arrows.meta,positioning,calc,intersections,angles,quotes}
\usepgfplotslibrary{fillbetween}
\pgfplotsset{compat=1.18}
\definecolor{primary}{RGB}{30,75,145}
\definecolor{primarydark}{RGB}{20,50,100}
\definecolor{defcol}{RGB}{0,114,178}
\definecolor{thmcol}{RGB}{0,135,110}
\definecolor{excol}{RGB}{0,130,85}
\definecolor{remarkcol}{RGB}{120,70,180}
\definecolor{attncol}{RGB}{200,40,55}
\definecolor{methcol}{RGB}{85,85,100}
\definecolor{areacol}{RGB}{120,160,230}
\definecolor{areacol2}{RGB}{230,150,80}
\newcommand{\dd}{\mathrm{d}}
\newcommand{\R}{\mathbb{R}}
\newcommand{\ua}{\,\text{u.a.}}
\newcommand{\tg}{\tan}\newcommand{\cotg}{\cot}
\tcbset{boxbase/.style={enhanced,breakable,boxrule=0.9pt,arc=2.5pt,left=3mm,right=3mm,top=2.3mm,bottom=2.3mm,fonttitle=\bfseries,coltitle=white,toptitle=0.6mm,bottomtitle=0.6mm}}
\newtcolorbox{methode}[1][]{boxbase,colback=methcol!7,colframe=methcol,sharp corners,title={\textsc{Comment faire}\ifx\relax#1\relax\relax\else\textmd{ : \itshape #1}\fi}}
\newtcolorbox{exemple}[1][]{boxbase,colback=excol!5,colframe=excol,title={\textsc{Exemple}\ifx\relax#1\relax\relax\else\textmd{ : \itshape #1}\fi}}
\newtcolorbox{idee}[1][]{boxbase,colback=defcol!5,colframe=defcol,title={\textsc{Idée clé}\ifx\relax#1\relax\relax\else\textmd{ : \itshape #1}\fi}}
\newtcolorbox{remarque}[1][]{boxbase,colback=remarkcol!6,colframe=remarkcol,title={\textsc{Remarque}\ifx\relax#1\relax\relax\else\textmd{ : \itshape #1}\fi}}
\newtcolorbox{attention}[1][]{boxbase,colback=attncol!6,colframe=attncol,title={\textsc{Attention}\ifx\relax#1\relax\relax\else\textmd{ : \itshape #1}\fi}}
\newtcolorbox{exo}[1][]{boxbase,colback={primary!4},colframe=primary,title={\textsc{Exercice #1}}}
\newtcolorbox{corr}[1][]{boxbase,colback={thmcol!5},colframe=thmcol,title={\textsc{Corrigé #1}}}
\tikzset{axe/.style={->,>=Stealth,black!65,line width=0.7pt},courbe/.style={primary,line width=1.3pt},courbe2/.style={areacol2!90!black,line width=1.3pt},dvert/.style={gray,dashed,line width=0.7pt}}
\pagestyle{fancy}\fancyhf{}
\fancyhead[L]{\small\textcolor{primarydark}{\textbf{Fiche 7} — Intégrale d'une fonction continue}}
\fancyhead[R]{\small\textcolor{primarydark}{\textsc{Thème 5 — Calculs d'aires}}}
\fancyfoot[C]{\small\thepage}
\renewcommand{\headrulewidth}{0.8pt}
\titleformat{\section}{\large\bfseries\color{primarydark}}{\thesection}{0.6em}{}

\begin{document}

\section*{Exercices — Facile : aire sous une courbe positive}

\begin{center}\itshape\small
Dans tous ces exercices, $f$ est positive sur l'intervalle considéré : l'aire cherchée est
directement égale à l'intégrale. Donner le résultat exact en u.a.
\end{center}

\begin{exo}[F1 — lecture directe]
Soit $f(x)=3x^2$, positive sur $\R$. Calculer l'aire $\mathcal{A}$ du domaine délimité par
$\mathcal{C}_f$, l'axe des abscisses et les droites $x=0$ et $x=2$.

\smallskip
\begin{center}
\begin{tikzpicture}
  \begin{axis}[width=7.4cm,height=4.4cm,axis lines=middle,
      xmin=-0.3,xmax=2.6,ymin=-1.2,ymax=13,xtick={0,2},ytick=\empty,
      xlabel={$x$},ylabel={$y$},
      every axis x label/.style={at={(current axis.right of origin)},anchor=west},
      every axis y label/.style={at={(current axis.above origin)},anchor=south}]
    \addplot[domain=-0.1:2.3,samples=60,courbe]{3*x^2};
    \addplot[name path=F,domain=0:2,samples=40,draw=none,forget plot]{3*x^2};
    \path[name path=b] (axis cs:0,0)--(axis cs:2,0);
    \addplot[areacol!45] fill between[of=F and b];
    \node at (axis cs:1.5,3){$\mathcal{A}$};
    \node[primary] at (axis cs:2.05,11.5){$f$};
  \end{axis}
\end{tikzpicture}
\end{center}
\end{exo}

\begin{exo}[F2 — fonction affine]
Soit $f(x)=x+1$. Calculer l'aire du domaine situé sous $\mathcal{C}_f$, au-dessus de l'axe des
abscisses, entre $x=0$ et $x=4$.
\end{exo}

\begin{exo}[F3 — bornes par $f(x)=0$]
Soit $f(x)=4-x^2$.
\begin{enumerate}[label=\alph*),leftmargin=2em,topsep=2pt]
  \item Résoudre $f(x)=0$ ; en déduire les deux bornes du domaine où $\mathcal{C}_f$ est au-dessus de l'axe.
  \item Vérifier que $f$ est positive entre ces deux bornes, puis calculer l'aire correspondante.
\end{enumerate}
\end{exo}

\begin{exo}[F4 — puissance sur un intervalle décalé]
Soit $f(x)=x^2$. Calculer l'aire du domaine délimité par $\mathcal{C}_f$, l'axe des abscisses et
les droites $x=1$ et $x=3$.
\end{exo}

\begin{exo}[F5 — bornes par $f(x)=0$, valeur symétrique]
Soit $f(x)=9-x^2$.
\begin{enumerate}[label=\alph*),leftmargin=2em,topsep=2pt]
  \item Déterminer les racines de $f$.
  \item Calculer l'aire de la portion de plan comprise entre $\mathcal{C}_f$ et l'axe des abscisses.
\end{enumerate}
\end{exo}

\end{document}
